\documentclass[11pt]{article}
\usepackage[margin=1in]{geometry}
\usepackage[T1]{fontenc}
\usepackage[utf8]{inputenc}
\usepackage{lmodern,amsmath,amssymb,amsthm,mathtools,booktabs,microtype}
\usepackage[colorlinks=true,linkcolor=blue,citecolor=blue,urlcolor=blue]{hyperref}
\hypersetup{pdftitle={Agnostic oracle bounds for Bayesian neural networks},pdfauthor={ChatGPT-6 Astra; Claude Opus 5.5; Kaarel Hänni; Dmitry Vaintrob}}
\newtheorem{theorem}{Theorem}[section]
\newtheorem{lemma}[theorem]{Lemma}
\newtheorem{proposition}[theorem]{Proposition}
\newtheorem{corollary}[theorem]{Corollary}
\theoremstyle{definition}
\newtheorem{definition}[theorem]{Definition}
\theoremstyle{remark}
\newtheorem{remark}[theorem]{Remark}
\newcommand{\E}{\mathbb E}
\newcommand{\R}{\mathbb R}
\newcommand{\N}{\mathbb N}
\newcommand{\Prob}{\mathbb P}
\newcommand{\clip}{\operatorname{clip}}
\newcommand{\KL}{\operatorname{KL}}
\newcommand{\Var}{\operatorname{Var}}
\newcommand{\Rad}{\operatorname{Rad}}
\newcommand{\Risk}{R}
\newcommand{\empRisk}{\widehat R}
\DeclareMathOperator*{\argmin}{arg\,min}
\title{Agnostic oracle bounds for Bayesian neural networks}
\author{ChatGPT-6 Astra \and Claude Opus 5.5 \and Kaarel H\"anni \and
 Dmitry Vaintrob\thanks{Authors are listed in alphabetical order.}}
\date{14 November 2026}
\begin{document}
\maketitle
\begin{abstract}
We study whether Bayesian neural-network prediction can compete with the best small network when the data need not be generated by any network. Labels are bounded, inputs range over a continuous domain, and training observations are independent and identically distributed. We give two guarantees with coefficient one on the comparator's population squared loss. First, for a fully connected Gaussian tanh prior in the strong-prior limit, we improve a tempered-likelihood oracle bound by retaining the learner's energy penalty in the statistical comparison. At fixed depth $L\ge3$, the remainder decreases as $\widetilde O(n^{-1/(L-1)})$, instead of $\widetilde O(n^{-1/(L+1)})$, at every learner width at least the teacher width; one hidden layer gives a square-root rate. Related bounds hold at finite Gaussian prior scale for every $\gamma\ge2$, and for every $\gamma\ge1$ with one hidden layer. Second, a prior that explicitly mixes over network architectures gives a final-posterior-mean bound in expectation of order the teacher parameter count divided by $n$, with additional logarithmic and polynomial-depth factors. This second result covers tanh, ReLU and GELU, but changes the prior. It uses a continuous-prior version of a recent exponential-weights aggregation theorem, for which we provide a short proof. We distinguish these results from exact interpolation, width-only comparison with arbitrarily large teacher weights, and fast high-probability guarantees for Bayesian model averaging.
\end{abstract}

\section{The approximation question}
Real data need not have a small exact neural implementation. A useful learning guarantee should instead compare a predictor with the best small network that approximates the data. Let $D$ be any distribution on $X\times[-1,1]$, with $X\subseteq[-1,1]^d$, and let
\[
 \Risk(f)=\E_{(x,y)\sim D}(f(x)-y)^2,
 \qquad \empRisk_S(f)=\frac1n\sum_{i=1}^n(f(x_i)-y_i)^2.
\]
We seek bounds of the form
\[
 \Risk(\widehat f_S)\le \Risk(g)+\text{a complexity penalty tending to zero},
\]
where $g$ is a small multilayer perceptron and the learner may have arbitrarily larger width. Such a bound is \emph{agnostic}: it makes no realizability assumption, and permits arbitrary bounded noise and model misspecification. An additive remainder is necessary even for estimating an unknown constant regression mean.

There are two different sources of a favorable prior. One can start with independent Gaussian weights in a fixed, possibly very wide architecture and prove that their induced preferences are useful. Alternatively, one can explicitly assign prior mass to small architectures. We treat both approaches, and keep their conclusions separate.

\paragraph{Gaussian weights and a tempered likelihood.}
For tanh networks with $L-1$ hidden layers, our first theorem uses exactly the Gaussian energy and raw-output likelihood defined in Section~\ref{sec:gaussian-oracle}. The prior scale tends to zero after the sample and learner width are fixed. A previous proof~\cite{opus} bounded the complexity of every empirical minimizer using only its maximum possible energy. We instead retain the negative learner-energy term when comparing with a teacher. Statistical deviations can then be absorbed by this term. The resulting bound holds with high probability on one event for every eligible learner width, for posterior draws in risk expectation and for the mean of clipped predictions. The raw-output extension in Appendix~\ref{sec:raw-gaussian-oracle} uses a different calibration and a standard smooth-loss theorem. Section~\ref{sec:finite-gaussian-oracle} also gives finite-prior versions, without a prior-strength limit. The depth-dependent sample exponent remains a substantial limitation.

\paragraph{An explicit prior over architectures.}
Our second theorem uses a fixed Gaussian observation likelihood centered on clipped predictions and a prior that chooses a network size, then continuous weights. It gives a bound for the final posterior mean using all observations, with no zero-temperature limit or averaging over training prefixes. At fixed input dimension and depth, a width-$m$ teacher with polynomially bounded weights has an expected excess bound $\widetilde O(m^2/n)$. Consequently the learner can compete with widths $m\le n^\alpha$ for every fixed $\alpha<1/2$. This is a stronger approximation guarantee, but it builds the preference for small architectures into the prior. It does not establish the same bound for independent Gaussian weights on every neuron of a large fixed architecture.

\paragraph{Relation to aggregation theory.}
H\o gsgaard, Rebeschini and Wegel~\cite{hrw} established the sharp expected-risk behavior of final exponential-weights aggregation under random design, resolving a question about this estimator at sufficiently high constant temperature. The statistical principle behind our second theorem belongs to this line of work. We give a self-contained continuous-prior argument, with a conservative temperature constant, and a quantitative neural-network application. Fast expectation bounds do not automatically imply fast bounds with high probability.

\section{A sharper oracle inequality for a Gaussian weight prior}
\label{sec:gaussian-oracle}

This section concerns a fully connected network with $H\geq1$ hidden
layers, all of width $N$, and a scalar linear readout.  Its total depth,
counting the readout, is $L=H+1$.  The activation in every hidden layer is
$\tanh$.  Writing $\bar x=(1,x)\in\mathbb R^{d+1}$, the network is
\[
 h^1(x)=\tanh(W^1\bar x),\qquad
 h^\ell(x)=\tanh(W^\ell h^{\ell-1}(x))\quad(2\leq\ell\leq H),
 \qquad z_w(x)=v^\top h^H(x).
\]
The activation is applied coordinatewise.  There are no internal biases;
the constant input coordinate supplies the first-layer biases.  Define
the weighted quadratic energy
\begin{equation}
 Q_N(w)=\frac dN\|W^1\|_F^2+
       \sum_{\ell=2}^H\|W^\ell\|_F^2+N\|v\|_2^2.
 \label{eq:gaussian-energy}
\end{equation}
Here and below $d\geq1$.  Inputs may have any distribution on
$\mathbb R^d$; the labels satisfy $|Y|\leq1$.  No bound on the input norm
is needed for the argument in this section.

Let $S=((X_i,Y_i))_{i=1}^n$ consist of independent samples from a law $D$.
For a predictor $g$, write
\[
 R(g)=\mathbb E_D(g(X)-Y)^2,\qquad
 \widehat R_S(g)=\frac1n\sum_{i=1}^n(g(X_i)-Y_i)^2,
 \qquad \operatorname{clip}(t)=\max\{-1,\min\{1,t\}\}.
\]
For positive $\sigma$ and $\lambda$, give the weights independent
centered Gaussian priors with variances
\begin{equation}
 \operatorname{Var}(W^1_{jk})=\frac{N\sigma^2}{d},\qquad
 \operatorname{Var}(W^\ell_{jk})=\sigma^2\ (2\leq\ell\leq H),\qquad
 \operatorname{Var}(v_j)=\frac{\sigma^2}{N}.
 \label{eq:gaussian-prior-variances}
\end{equation}
Use a Gaussian likelihood with mean $z_w(X_i)$ and variance
$n\lambda\sigma^2$ for each observation.  The posterior is therefore
\begin{equation}
 \Pi_{\sigma,N,S}(dw)=\frac1{Z_{\sigma,N,S}}
  \exp\left[-\frac{Q_N(w)+\lambda^{-1}\widehat R_S(z_w)}{2\sigma^2}\right]dw.
 \label{eq:gaussian-posterior}
\end{equation}
The likelihood uses the raw real-valued network output; it does not
require exact interpolation.  The limit $\sigma\downarrow0$ is taken
with $n,S,N$, and $\lambda$ fixed.  Its limiting optimization problem is
\begin{equation}
 \operatorname*{minimize}_w
 \bigl\{\widehat R_S(z_w)+\lambda Q_N(w)\bigr\}.
 \label{eq:gaussian-objective}
\end{equation}

\subsection{Teacher implementations and a width-independent capacity bound}

\begin{lemma}[Cloning a teacher]
\label{lem:gaussian-cloning}
Suppose that $g$ is a network of the preceding architecture, hidden
width $m$, and depth $H+1$, with every weight of absolute value at most
$B$.  At every width $N\geq m$ it has an exact implementation with
\begin{equation}
 Q_N\leq E:=d(d+1)B^2+2H m^2B^2.
 \label{eq:gaussian-teacher-energy}
\end{equation}
Moreover, $|g(x)|\leq U:=mB$ for every $x$.
\end{lemma}
\begin{proof}
Partition the $N$ neurons in each hidden layer into $m$ groups, of sizes
$k_j\in\{\lfloor N/m\rfloor,\lceil N/m\rceil\}$.  Copy a teacher
first-layer row onto every neuron in its group.  Replace a teacher
internal weight $A^\ell_{ji}$ by $A^\ell_{ji}/k_i$ on every corresponding
edge, and its readout weight $u_j$ by $u_j/k_j$ on each copy.  Every copy
has exactly the teacher activation, so the function is unchanged.
The first-layer energy is at most $d(d+1)B^2$.  Since $k_j/k_i\leq2$,
each internal layer costs at most $2m^2B^2$.  Since
$k_j\geq N/(2m)$, the readout costs at most $2m^2B^2$.  There are
$H-1$ internal layers and one readout.  Finally, $|\tanh|\leq1$ gives
$|g|\leq\sum_j|u_j|\leq mB$.
\end{proof}

Put
\begin{equation}
 D_n=\sqrt{2(d+1)\log(n+1)}+\frac{H-1}{2},\qquad
 \mathcal Z_K=\bigcup_{N\geq1}\{z_w:Q_N(w)\leq K\}.
 \label{eq:gaussian-Dn}
\end{equation}
The union over all learner widths is essential: the following event
controls every width at once.

\begin{lemma}[Capacity of an energy ball]
\label{lem:gaussian-capacity}
For every set of $n$ inputs and every $K>0$, the empirical Rademacher
complexity satisfies
\begin{equation}
 \widehat{\mathfrak R}_S(\mathcal Z_K)
 \leq \left(\frac KH\right)^{H/2}\frac{D_n}{\sqrt n}.
 \label{eq:gaussian-rademacher}
\end{equation}
Consequently, with probability at least $1-\eta$, simultaneously for all
widths and all $w$ with $Q_N(w)\leq K$,
\begin{equation}
 R(\operatorname{clip}z_w)
 \leq\widehat R_S(\operatorname{clip}z_w)
 +8\left(\frac KH\right)^{H/2}\frac{D_n}{\sqrt n}
 +4\sqrt{\frac{\log(1/\eta)}{2n}}.
 \label{eq:gaussian-fixed-energy-deviation}
\end{equation}
\end{lemma}
\begin{proof}
For a bounded set $A\subset\mathbb R^n$, write
$\operatorname{Rad}(A)=n^{-1}\mathbb E_\epsilon
\sup_{a\in A}\sum_i\epsilon_i a_i$.
Its value is unchanged by taking a closed convex hull; it scales
linearly under positive scalar multiplication.  We use the standard
Rademacher contraction inequality, Massart's finite-class bound, and
the Sauer--Shelah bound for halfspaces.

For $t\in[0,1]$, let $\operatorname{clip}_t(u)=\max\{-t,\min\{t,u\}\}$,
and let $\operatorname{Clip}(A)$ contain the vectors
$(\operatorname{clip}_t(a_i))_{i=1}^n$ for $a\in A$ and a single common
$t\in[0,1]$.  The identity
$\operatorname{clip}_t(u)=(|u+t|-|u-t|)/2$, contraction, and symmetry of
the Rademacher signs imply
\begin{equation}
 \operatorname{Rad}(\operatorname{Clip}(A))
 \leq\operatorname{Rad}(A)+\frac1{2\sqrt n}.
 \label{eq:gaussian-variable-clip}
\end{equation}
Indeed, split the supremum between the two absolute-value expressions
and contract each.  The two additional terms sum to
$\mathbb E|\sum_i\epsilon_i|/(2n)\leq1/(2\sqrt n)$.

We also need a positive mixture identity.  For $s>0$ and $|u|\leq1$,
\begin{equation}
 \frac{\tanh(su)}s=
 \int_0^1\operatorname{clip}_t(u)\,
       2s\tanh(st)\operatorname{sech}^2(st)\,dt
 +\operatorname{sech}^2(s)\operatorname{clip}_1(u).
 \label{eq:gaussian-tanh-mixture}
\end{equation}
All weights are nonnegative and have total mass one.  To verify the
identity for $u\geq0$, integrate by parts with
$\psi(t)=\tanh^2(st)$: the right side is
$u-\int_0^u\tanh^2(st)dt=\tanh(su)/s$.  Oddness gives the other half.

On the fixed sample inputs, let $\mathcal F_1$ be the closed convex hull
of the vectors $(\tanh(a\cdot\bar X_i))_{i=1}^n$, with
$a\in\mathbb R^{d+1}$.  Recursively set
$\mathcal F_\ell=\overline{\operatorname{conv}}
\operatorname{Clip}(\mathcal F_{\ell-1})$.  These sets are symmetric,
convex, contain zero, and lie in $[-1,1]^n$.
Define $\rho^1=\boldsymbol1_N$ and
$\rho^\ell=|W^\ell|\rho^{\ell-1}$, with entrywise absolute values.
Induction using \eqref{eq:gaussian-tanh-mixture} shows that the sample
vector of hidden neuron $j$ at layer $\ell$ belongs to
$\rho^\ell_j\mathcal F_\ell$.  A zero value of $\rho^\ell_j$ forces the
neuron to vanish.  For a positive value, divide the preactivation by
$\rho^\ell_j$; it is a convex signed combination of the previous
normalized neurons, hence belongs to $\mathcal F_{\ell-1}$.  The mixture
identity completes the induction.

It follows that the output sample vector belongs to
$P(w)\mathcal F_H$, where
$P(w)=|v|^\top|W^H|\cdots|W^2|\boldsymbol1_N$ (with no matrix factors
when $H=1$).  Cauchy--Schwarz and the Frobenius bound on operator norm give
\[
 P(w)\leq\sqrt N\|v\|_2\prod_{\ell=2}^H\|W^\ell\|_F
 \leq\left(\frac{Q_N(w)}H\right)^{H/2}.
\]
The last inequality is AM--GM over the $H$ energy terms after the first
layer.  The first-layer energy is not needed for this capacity estimate.
By \eqref{eq:gaussian-variable-clip},
$\operatorname{Rad}(\mathcal F_H)\leq
\operatorname{Rad}(\mathcal F_1)+(H-1)/(2\sqrt n)$.

Finally, if a scalar random variable $T$ has distribution function
$(1+\tanh t)/2$, then $\tanh u=\mathbb E\operatorname{sign}(u-T)$.
Thus $\mathcal F_1$ lies in the convex hull of the affine-halfspace
dichotomies of the sample inputs.  There are at most $(n+1)^{d+1}$ such
dichotomies by Sauer--Shelah.  Massart's bound gives
$\operatorname{Rad}(\mathcal F_1)
\leq\sqrt{2(d+1)\log(n+1)/n}$, proving
\eqref{eq:gaussian-rademacher}.

The clipped squared loss is bounded by $4$ and is $4$-Lipschitz in the
raw prediction.  Symmetrization and contraction bound its expected
uniform deviation by eight times
\eqref{eq:gaussian-rademacher}.  Changing one observation changes that
supremum by at most $4/n$, so bounded-difference concentration gives
\eqref{eq:gaussian-fixed-energy-deviation}.
The suprema are measurable: for each width the relevant loss is
continuous in the parameters, and rational parameter vectors in the
open energy ellipsoid are dense in the closed ellipsoid.  Taking the
union over the countably many widths gives a countable determining
class.
\end{proof}

\subsection{Absorbing capacity into the energy penalty}

An estimate based only on the largest possible energy of a minimizer
loses useful information.  In the proof below, the capacity cost of
the selected network is paid by that same network's energy penalty.

\begin{theorem}[Gaussian oracle inequality, at least two hidden layers]
\label{thm:gaussian-absorption}
Suppose $H\geq2$, $0<\delta<1$, and $n\geq1$.  Set
\begin{align}
 a_n&=8(2/H)^{H/2}\frac{D_n}{\sqrt n},&
 \lambda&=a_n^{2/H}
       =\frac2H\left(\frac{8D_n}{\sqrt n}\right)^{2/H},
       \label{eq:gaussian-calibration}\\
 J&=\max\{0,\lceil\log_2(1/\lambda)\rceil\},&
 b_n&=4\sqrt{\frac{\log(2(J+1)/\delta)}{2n}},&
 e_n&=8H^{-H/2}\frac{D_n}{\sqrt n}.
 \label{eq:gaussian-residuals}
\end{align}
Fix a teacher $g$ as in Lemma~\ref{lem:gaussian-cloning}, independently
of the observed sample.  With probability at least $1-\delta$,
simultaneously for every learner width $N\geq m$, every global
minimizer $w$ of \eqref{eq:gaussian-objective} satisfies
\begin{equation}
 R(\operatorname{clip}z_w)
 \leq R(g)+\lambda E+e_n+b_n+h_g,
 \label{eq:gaussian-oracle}
\end{equation}
where one may use either of the following teacher concentration bounds,
chosen in advance (each choice gives its own event of probability at
least $1-\delta$):
\begin{align}
 h_g&=(U+1)^2\sqrt{\frac{\log(2/\delta)}{2n}},
       \label{eq:gaussian-teacher-hoeffding}\\
 h_g&=(U+1)\sqrt{\frac{2R(g)\log(2/\delta)}n}
       +\frac{2(U+1)^2\log(2/\delta)}{3n}.
       \label{eq:gaussian-teacher-bernstein}
\end{align}
On that same event, for each fixed $N\geq m$, the right side of
\eqref{eq:gaussian-oracle} also bounds
\[
 \limsup_{\sigma\downarrow0}
   \mathbb E_{w\sim\Pi_{\sigma,N,S}}R(\operatorname{clip}z_w)
 \quad\hbox{and}\quad
 \limsup_{\sigma\downarrow0}
   R\left(\mathbb E_{w\sim\Pi_{\sigma,N,S}}\operatorname{clip}z_w\right).
\]
The event is common to all widths.  The limit in $\sigma$ is pointwise
in the learner width, with no width-uniform convergence rate asserted.
\end{theorem}
\begin{proof}
The zero network has empirical risk at most one.  Consequently every
global minimizer has energy $q=Q_N(w)\leq1/\lambda$.
Apply Lemma~\ref{lem:gaussian-capacity} at thresholds
$K_j=2^j$, $j=0,\ldots,J$, each with failure probability
$\delta/[2(J+1)]$.  On the resulting event, common to every width,
networks with $q\leq1$ have deviation at most $e_n+b_n$.  For $q>1$,
the least dyadic threshold containing $q$ is less than $2q$, so their
deviation is at most $a_n q^{H/2}+b_n$.  In particular, for every
$q\leq1/\lambda$,
\[
 R(\operatorname{clip}z_w)
 \leq\widehat R_S(\operatorname{clip}z_w)
       +a_nq^{H/2}+e_n+b_n.
\]
Since $H/2\geq1$ and $\lambda^{H/2}=a_n$,
\[
 a_nq^{H/2}\leq a_n\lambda^{1-H/2}q=\lambda q.
\]
Clipping cannot increase squared loss against labels in $[-1,1]$.
Comparison of the minimizing objective with the cloned teacher gives
\[
 R(\operatorname{clip}z_w)
 \leq\widehat R_S(z_w)+\lambda Q_N(w)+e_n+b_n
 \leq\widehat R_S(g)+\lambda E+e_n+b_n.
\]
Hoeffding's inequality gives
\eqref{eq:gaussian-teacher-hoeffding} with failure probability at most
$\delta/2$.  For \eqref{eq:gaussian-teacher-bernstein}, put
$Z=(g(X)-Y)^2$.  Then $0\leq Z\leq(U+1)^2$ and
$\operatorname{Var}(Z)\leq(U+1)^2R(g)$; Bernstein's inequality gives
the stated bound.  Union-bounding with the capacity event proves
\eqref{eq:gaussian-oracle}.

For completeness, at fixed $N,S$ the function
$F(w)=Q_N(w)+\lambda^{-1}\widehat R_S(z_w)$ is continuous and coercive,
so its minimizer set is nonempty and compact.  Write $f_* =\min F$.
For each $\varepsilon>0$, a positive-volume neighborhood of a minimizer
has $F\leq f_*+\varepsilon/2$.  If $\beta=(2\sigma^2)^{-1}$ and
$\beta>\beta_0>0$, then
\[
 \int_{\{F\geq f_*+\varepsilon\}}e^{-\beta F}dw
 \leq e^{-(\beta-\beta_0)(f_*+\varepsilon)}
       \int e^{-\beta_0F}dw.
\]
The final integral is finite because $F\geq Q_N$.  Dividing by the
normalizer lower bound from that neighborhood shows that the posterior
mass of $\{F\geq f_*+\varepsilon\}$ tends to zero.  Coercivity and
continuity then show that the posterior expectation of any bounded
continuous function has upper limit at most its maximum on the
minimizer set.  Apply this to $R(\operatorname{clip}z_w)\in[0,4]$;
its continuity follows by dominated convergence.  Jensen's inequality
gives the assertion for the mean of clipped outputs.
\end{proof}

The leading teacher-dependent term is explicitly
\begin{equation}
 \lambda E=
 \left[\frac{2d(d+1)B^2}{H}+4m^2B^2\right]
 \left(\frac{8D_n}{\sqrt n}\right)^{2/H}.
 \label{eq:gaussian-leading-term}
\end{equation}
Thus, at fixed depth and for a fixed teacher, the excess risk is
$\widetilde O(n^{-1/H})=\widetilde O(n^{-1/(L-1)})$.
For two hidden layers it is $\widetilde O(n^{-1/2})$.
The likelihood variance is $n\lambda\sigma^2$ with the exact
$\lambda$ in \eqref{eq:gaussian-calibration}; it is not a sample-independent
constant.  In particular, up to logarithms this calibration lies at the
power $\lambda\asymp n^{-1/H}$.

\begin{theorem}[One hidden layer]
\label{thm:gaussian-one-hidden}
Suppose $H=1$, and put
\[
 a_n=8\sqrt2\,D_n/\sqrt n,\qquad e_n=8D_n/\sqrt n.
\]
For any $\lambda>0$, define $J,b_n$ by
\eqref{eq:gaussian-residuals}.  With the same assumptions and posterior
interpretation as Theorem~\ref{thm:gaussian-absorption}, the bound is
\begin{equation}
 R(\operatorname{clip}z_w)
 \leq R(g)+\lambda E+\frac{a_n^2}{4\lambda}+e_n+b_n+h_g.
 \label{eq:gaussian-one-hidden-oracle}
\end{equation}
In particular, $\lambda=a_n$ gives an excess of order
$(E+1)D_n/\sqrt n$ plus the concentration terms.  If a deterministic
energy cap $E_0\geq E$ is known, choosing
\begin{equation}
 \lambda=\frac{a_n}{2\sqrt{E_0+1}}
 \label{eq:gaussian-one-hidden-calibration}
\end{equation}
replaces the two energy terms in
\eqref{eq:gaussian-one-hidden-oracle} by at most
$a_n\sqrt{E_0+1}$.
\end{theorem}
\begin{proof}
The same dyadic argument gives deviation at most
$a_n\sqrt q+e_n+b_n$.  After comparing penalized empirical risks, use
\[
 \sup_{q\geq0}\{a_n\sqrt q-\lambda q\}=\frac{a_n^2}{4\lambda}.
\]
The teacher concentration and posterior-limit arguments are unchanged.
For the stated calibration,
$\lambda E+a_n^2/(4\lambda)
=a_n[E/(2\sqrt{E_0+1})+\sqrt{E_0+1}/2]
\leq a_n\sqrt{E_0+1}$.
\end{proof}

\subsection{Comparison with an expanding class of teachers}

\begin{corollary}[Growing teacher width and weight bound]
\label{cor:gaussian-growing-teachers}
Fix $H,d$ and $0<\delta<1$.  Let $\mathcal G_n$ be the class of networks
with $H$ hidden layers, width at most $\lfloor n^\alpha\rfloor$, and
every weight bounded in absolute value by $n^\alpha$.

If $H\geq2$, use \eqref{eq:gaussian-calibration}.  For every
$0\leq\alpha<1/(4H)$, with probability at least $1-\delta$,
simultaneously for all learner widths $N\geq\lfloor n^\alpha\rfloor$,
\begin{equation}
 \limsup_{\sigma\downarrow0}
  \mathbb E_{\Pi_{\sigma,N,S}}R(\operatorname{clip}z_w)
 \leq\min_{g\in\mathcal G_n}R(g)
       +\widetilde O_{H,d,\delta}(n^{-(1/H-4\alpha)}).
 \label{eq:gaussian-growing-deep}
\end{equation}
If $H=1$, use \eqref{eq:gaussian-one-hidden-calibration} with
$E_0=[d(d+1)+2]n^{4\alpha}$.  For every
$0\leq\alpha<1/4$, the corresponding bound is
\begin{equation}
 \limsup_{\sigma\downarrow0}
  \mathbb E_{\Pi_{\sigma,N,S}}R(\operatorname{clip}z_w)
 \leq\min_{g\in\mathcal G_n}R(g)
       +\widetilde O_{d,\delta}(n^{-(1/2-2\alpha)}).
 \label{eq:gaussian-growing-shallow}
\end{equation}
Both assertions also hold for the risk of the posterior mean of clipped
outputs.  The notation $\widetilde O$ suppresses logarithms in $n$;
the preceding theorems give explicit finite-sample constants.
\end{corollary}
\begin{proof}
Narrower teachers embed by zero padding, since $\tanh(0)=0$.
The parameter set defining $\mathcal G_n$ is compact and its population
risk is continuous, so it has a deterministic population-risk minimizer
$g_n$.  The zero network is an available comparator, hence $R(g_n)\leq1$.
Lemma~\ref{lem:gaussian-cloning} gives
\[
 E\leq[d(d+1)+2H]n^{4\alpha},\qquad U\leq n^{2\alpha}.
\]
The Bernstein term \eqref{eq:gaussian-teacher-bernstein} is therefore
$O_\delta(n^{2\alpha-1/2}+n^{4\alpha-1})$.
For $H\geq2$, \eqref{eq:gaussian-leading-term} has order
$\widetilde O_{H,d}(n^{4\alpha-1/H})$, which dominates the preceding
concentration powers.  For $H=1$, the energy-calibrated bound in
Theorem~\ref{thm:gaussian-one-hidden} has order
$\widetilde O_d(n^{2\alpha-1/2})$.  The dyadic concentration terms are
at most logarithmic factors times $n^{-1/2}$ in both cases.
\end{proof}

\paragraph{What the exponent does and does not establish.}
The argument is independent of learner width and allows arbitrary
label noise and arbitrary approximation error.  It still uses a weight
bound, or another bound on the teacher's energy, as well as its width.
It does not give a bound in teacher width alone for arbitrary
real-valued networks.  It also does not make the dependence on depth
polynomial: the leading sample requirement from
\eqref{eq:gaussian-leading-term} raises the teacher-energy and accuracy
factors to the power $H$.

The distinction between clipping before and after posterior averaging
is substantive.  The bounded-loss argument above controls
$\mathbb E_\Pi R(\operatorname{clip}z_w)$ and
$R(\mathbb E_\Pi\operatorname{clip}z_w)$.  By itself it does not control
$R(\mathbb E_\Pi z_w)$ or
$R(\operatorname{clip}(\mathbb E_\Pi z_w))$.

\section{A finite-scale Gaussian result with no prior-strength limit}
\label{sec:finite-gaussian-oracle}

A related oracle inequality holds for an actual finite-scale posterior,
without sending a prior parameter to zero.  The proof below applies to
all learner widths at once when the prior variances decay sufficiently
quickly with width.  It controls clipped posterior draws and the mean
of clipped outputs.  Its likelihood still uses the raw real-valued
network output.

Use the network architecture, energy $Q_N$, and risks from
Section~\ref{sec:gaussian-oracle}.  In this section inputs lie in
$[-1,1]^d$ and labels lie in $[-1,1]$; their joint law is otherwise
arbitrary.  For $\gamma\geq2$---and also for $\gamma\geq1$ in the
one-hidden-layer result below---let the prior $\mu_{N,\gamma}$ have
independent centered Gaussian coordinates with variances
\begin{gather}
 \operatorname{Var}(W^1_{jk})=\frac{N^{1-\gamma}}d,\qquad
 \operatorname{Var}(v_j)=N^{-\gamma-1},\notag\\
 \operatorname{Var}(W^\ell_{jk})=N^{-\gamma}\quad(2\leq\ell\leq H).
 \label{eq:finite-prior}
\end{gather}
Thus this is the Gaussian family of~\eqref{eq:gaussian-prior-variances}
with $\sigma^2=N^{-\gamma}$.  For $\lambda>0$, use a Gaussian likelihood
with variance $n\lambda N^{-\gamma}$ and mean $z_w(X_i)$.  The posterior
is
\begin{equation}
 \Pi_{N,\gamma,S}(dw)
   =\frac{\exp[-\beta_N\widehat R_S(z_w)]}{Z_{N,\gamma,S}}
          \mu_{N,\gamma}(dw),
 \qquad \beta_N=\frac{N^\gamma}{2\lambda}.
 \label{eq:finite-posterior}
\end{equation}
All statements below concern this probability measure at its specified
finite width and scale.

\subsection{An explicit finite-scale oracle inequality}

Let $g$ be a teacher of hidden width $m$, with weights bounded by
$B\geq1$.  Define
\begin{gather}
 E=d(d+1)B^2+2Hm^2B^2,\qquad U=mB,\notag\\
 D_0=d+H+1,\qquad A_g=\sqrt2mB,\notag\\
 K_g=1+A_g\left(\sqrt2 A_g^{H-1}
                       +\sum_{j=0}^{H-2}A_g^j\right).
 \label{eq:finite-teacher-constants}
\end{gather}
The sum is empty if $H=1$.  In particular, $K_g\geq1$ and
$\log K_g=O(H\log(2+mB))$.
For $H\geq2$, put
\begin{gather}
 p=H/2,\qquad
 A_n=\frac{8D_n}{H^p\sqrt n},\qquad
 J_\lambda=\left\lceil\log_2\frac8\lambda\right\rceil,
 \notag\\
 t_\lambda=\log\frac{2(J_\lambda+1)}\delta,
 \qquad
 r_{n,\lambda}=A_n+4\sqrt{\frac{t_\lambda}{2n}}.
 \label{eq:finite-statistical-constants}
\end{gather}
For the scalar teacher concentration term, write
\begin{equation}
 b_n(g,\delta)
   =(U+1)\sqrt{\frac{2R(g)\log(2/\delta)}n}
       +\frac{2(U+1)^2\log(2/\delta)}{3n}.
 \label{eq:finite-teacher-deviation}
\end{equation}

\begin{theorem}[A finite-scale oracle inequality, uniform over widths]
\label{thm:finite-gaussian-oracle}
Let $H\geq2$, $n\geq2$, $0<\delta<1$, and suppose
\begin{equation}
 \frac1n\leq\lambda\leq1,
 \qquad \lambda^p\geq2^{4p-2}A_n.
 \label{eq:finite-calibration}
\end{equation}
For every fixed teacher $g$, with probability at least $1-\delta$,
simultaneously for every $N\geq m$ and every $\gamma\geq2$,
\begin{align}
 \int R(\operatorname{clip}z_w)\,\Pi_{N,\gamma,S}(dw)
 &\leq R(g)+b_n(g,\delta)+r_{n,\lambda}\notag\\
 &\quad+\lambda\left[
 E+1+D_0\log2
       +2D_0\log\frac{4K_g(U+1)}\lambda\right].
 \label{eq:finite-oracle}
\end{align}
The same right-hand side bounds the risk of the posterior mean of
clipped outputs,
\[
 x\longmapsto\int\operatorname{clip}(z_w(x))\,
\Pi_{N,\gamma,S}(dw).
\]

For example, take
\begin{equation}
 \lambda_n=\max\left\{\frac1n,
                 \bigl(2^{4p-2}A_n\bigr)^{1/p}\right\},
 \label{eq:finite-schedule}
\end{equation}
for sufficiently large $n$ that $\lambda_n\leq1$.  At fixed
$d,H,\delta$ and teacher, the excess in~\eqref{eq:finite-oracle} is
$\widetilde O(n^{-1/H})$.  The bound has no learner-width factor and
requires no limit in either $N$ or the prior scale.
\end{theorem}

\subsection{A deviation inequality valid at every energy}

The first step strengthens the deviation bound used for the
strong-prior minimizers: here it must hold over the full parameter
space because the finite-scale posterior has unbounded support.

\begin{lemma}[Global clipped deviation]
\label{lem:finite-global-deviation}
Under~\eqref{eq:finite-calibration}, with probability at least
$1-\delta/2$, simultaneously for every network at every width,
\begin{equation}
 R(\operatorname{clip}z_w)
     \leq\widehat R_S(z_w)+\frac\lambda2 Q_N(w)+r_{n,\lambda}.
 \label{eq:finite-global-deviation}
\end{equation}
\end{lemma}
\begin{proof}
Apply~\eqref{eq:gaussian-fixed-energy-deviation} to the deterministic
classes $\mathcal Z_{2^j}$, $0\leq j\leq J_\lambda$, with individual
failure probability $\delta/[2(J_\lambda+1)]$.
On the intersection event, write $k=Q_N(w)$.  For $k\leq1$, the
unit-energy class gives the claim, since clipping cannot increase
empirical square loss.  For $1<k\leq8/\lambda$, the appropriate shell
has radius at most $2k$, so its capacity term is at most
\[
 A_n2^p k^p
 \leq A_n2^p(8/\lambda)^{p-1}k
 \leq\frac\lambda2 k,
\]
where the last inequality is exactly
$\lambda^p\geq2^{4p-2}A_n$.
Finally, for $k>8/\lambda$, the clipped population loss is at most
four and the empirical loss is nonnegative; therefore
\[
 R(\operatorname{clip}z_w)-\widehat R_S(z_w)
   \leq4<\frac\lambda2 k.
\]
This also covers arbitrarily large energies.
\end{proof}

\subsection{The posterior's variational comparison}

There are
\[
 P_N=N(d+2)+(H-1)N^2
\]
parameters.  Under the Gaussian prior, $N^\gamma Q_N$ has a
chi-square distribution with $P_N$ degrees of freedom.  Thus
\[
 \log\int\exp\left(\frac{N^\gamma Q_N}4\right)d\mu_{N,\gamma}
       =\frac{P_N}2\log2.
\]
The entropy variational inequality implies, for every comparison
measure $Q$,
\begin{equation}
 \operatorname{KL}(Q\Vert\mu_{N,\gamma})
   \geq\frac{N^\gamma}4\int Q_N\,dQ-\frac{P_N}2\log2.
 \label{eq:finite-kl-energy}
\end{equation}
The posterior~\eqref{eq:finite-posterior} minimizes
\[
 Q\longmapsto\int\widehat R_S(z_w)\,dQ
                   +2\lambda N^{-\gamma}
                          \operatorname{KL}(Q\Vert\mu_{N,\gamma}).
\]
Combining its variational optimality with~\eqref{eq:finite-kl-energy}
gives, for every comparison measure $Q$ of finite relative entropy,
\begin{align}
 &\int\left[\widehat R_S(z_w)+\frac\lambda2 Q_N(w)\right]
                          \Pi_{N,\gamma,S}(dw)\notag\\
 &\hspace{8mm}\leq
      \int\widehat R_S(z_w)\,Q(dw)
       +2\lambda N^{-\gamma}\operatorname{KL}(Q\Vert\mu_{N,\gamma})
       +\lambda\frac{P_N}{N^\gamma}\log2.
 \label{eq:finite-variational-comparison}
\end{align}
All quantities are finite for the Gaussian comparison measure used
below.  They are also finite for the posterior: its density relative
to the prior is bounded, and the prior has finite energy moments.
For $\gamma\geq2$ and $N\geq1$,
\begin{equation}
 \frac{P_N}{N^\gamma}\leq d+2+H-1=D_0.
 \label{eq:finite-dimension-ratio}
\end{equation}
For $H\geq2$, this is the point at which the condition
$\gamma\geq2$ enters.  When $H=1$, the quadratic term in $P_N$ is
absent: $P_N/N^\gamma\leq d+2=D_0$ already holds for every
$\gamma\geq1$.

\subsection{A Gaussian comparison measure near the teacher}

Let $w_*$ be the exact teacher clone from
Lemma~\ref{lem:gaussian-cloning}.  Denote the prior covariance by
$\Sigma_N$.  For $0<r\leq1$, define
\[
 Q_r=\mathcal N(w_*,r^2\Sigma_N).
\]
The elementary Gaussian relative-entropy formula gives
\begin{equation}
 \operatorname{KL}(Q_r\Vert\mu_{N,\gamma})
  =\frac{N^\gamma Q_N(w_*)}2
       +\frac{P_N}2\bigl(r^2-1-\log r^2\bigr)
  \leq\frac{N^\gamma E}2+P_N\log(1/r).
 \label{eq:finite-comparison-kl}
\end{equation}

We now bound the comparison measure's prediction error.  The clone's
internal matrices have operator norm at most $A_g=\sqrt2mB$, and
$\sqrt N\|v_*\|_2\leq A_g$.  For the internal matrices, identify each
group of clones with its normalized indicator vector.  The induced
matrix has entries $T_{ji}\sqrt{k_j/k_i}$, whose Frobenius norm is at
most $\sqrt2\,\|T\|_F\leq\sqrt2mB$; the readout bound follows directly
from the cloning construction.
For fixed $x\in[-1,1]^d$, put
\[
 e_\ell(x)
   =\left(\frac1N\mathbb E_{Q_r}
        \|h^\ell(x)-h_*^\ell(x)\|_2^2\right)^{1/2}.
\]
The activation is $1$-Lipschitz.  In the first layer,
\begin{equation}
 e_1(x)
  \leq rN^{(1-\gamma)/2}\frac{\|(1,x)\|_2}{\sqrt d}
  \leq\sqrt2\,r.
 \label{eq:finite-first-perturbation}
\end{equation}
For later hidden layers, use
\[
 W^\ell h^{\ell-1}-W_*^\ell h_*^{\ell-1}
   =W_*^\ell(h^{\ell-1}-h_*^{\ell-1})
                   +(W^\ell-W_*^\ell)h^{\ell-1}.
\]
The noise matrix in the second term is independent of
$h^{\ell-1}$, and $\|h^{\ell-1}\|_2^2\leq N$.  Hence Minkowski's
inequality and its Gaussian second moment give
\begin{equation}
 e_\ell(x)
   \leq A_g e_{\ell-1}(x)+rN^{(1-\gamma)/2}
   \leq A_g e_{\ell-1}(x)+r.
 \label{eq:finite-hidden-perturbation}
\end{equation}
Readout noise is likewise independent of the last hidden layer and
has prediction variance at most $r^2N^{-\gamma}$.  Consequently
\begin{align}
 \left(\mathbb E_{Q_r}|z_w(x)-g(x)|^2\right)^{1/2}
 &\leq A_g e_H(x)+rN^{-\gamma/2}\notag\\
 &\leq r\left[1+A_g\left(\sqrt2 A_g^{H-1}
                       +\sum_{j=0}^{H-2}A_g^j\right)\right]
   =rK_g.
 \label{eq:finite-output-perturbation}
\end{align}
This is a uniform \emph{pointwise second-moment} bound.  It does not
assert that one draw is uniformly close to the teacher over all inputs;
that stronger assertion is unnecessary here.

Choose
\begin{equation}
 \rho=\frac\lambda{4(U+1)},\qquad r=\frac\rho{K_g}.
 \label{eq:finite-comparison-radius}
\end{equation}
These satisfy $0<r\leq1$.  Minkowski's inequality in the product of
$Q_r$ and the empirical input measure gives
\begin{align}
 \int\widehat R_S(z_w)\,Q_r(dw)
 &\leq\widehat R_S(g)+2\rho\sqrt{\widehat R_S(g)}+\rho^2
 \notag\\
 &\leq\widehat R_S(g)+2\rho(U+1)+\rho^2
 \leq\widehat R_S(g)+\lambda.
 \label{eq:finite-comparison-risk}
\end{align}

\begin{proof}[Proof of Theorem~\ref{thm:finite-gaussian-oracle}]
Integrate the global deviation inequality
\eqref{eq:finite-global-deviation} against the posterior, and then
apply~\eqref{eq:finite-variational-comparison} with $Q=Q_r$.
Equations~\eqref{eq:finite-dimension-ratio},
\eqref{eq:finite-comparison-kl}, and
\eqref{eq:finite-comparison-risk} imply
\begin{align*}
 \int R(\operatorname{clip}z_w)\,\Pi_{N,\gamma,S}(dw)
 &\leq\widehat R_S(g)+r_{n,\lambda}\\
 &\quad+\lambda\left[E+1+D_0\log2
               +2D_0\log\frac{4K_g(U+1)}\lambda\right].
\end{align*}
For the fixed teacher, Bernstein's inequality gives
$\widehat R_S(g)\leq R(g)+b_n(g,\delta)$ except on an event of
probability $\delta/2$.  Intersecting this event with the global
deviation event proves~\eqref{eq:finite-oracle}.  Neither event depends
on $N$ or $\gamma$, so the conclusion is simultaneous over all the
claimed widths and scales.  Jensen's inequality proves the claim for
the mean of clipped outputs.  The rate follows from
\eqref{eq:finite-schedule}; $\log(1/\lambda_n)$ adds a logarithm,
without changing its power of $n$.
\end{proof}

\subsection{One hidden layer and the limits of this argument}

At $H=1$, the finite-scale result holds for \emph{every}
$\gamma\geq1$, including the mean-field scale $\gamma=1$, with rate
$\widetilde O(n^{-1/2})$.  Its guarantee is simultaneous over all
$N\geq m$ and all these $\gamma$.  The wider range follows because
$P_N=(d+2)N$, so $P_N/N^\gamma\leq d+2$ for $\gamma\geq1$.
The Gaussian perturbation estimates also hold throughout this range:
$N^{(1-\gamma)/2}\leq1$ and $N^{-\gamma/2}\leq1$ are their only
scale-dependent requirements.  In
\eqref{eq:finite-statistical-constants} take $p=1/2$ and replace the
residual by
\[
 r_{n,\lambda}=A_n+\frac{A_n^2}{\lambda}
                        +4\sqrt{\frac{t_\lambda}{2n}}.
\]
There is no condition $\lambda^p\geq2^{4p-2}A_n$ in this case.
The needed shell estimate is simply
\[
 A_n\sqrt{2k}\leq\frac\lambda2 k+\frac{A_n^2}{\lambda}.
\]
The unit-energy and high-energy cases are unchanged, and all subsequent
steps of the proof hold verbatim.  Choosing
$\lambda_n=\max\{A_n,1/n\}$, for sufficiently large $n$ that
$\lambda_n\leq1$, gives the stated rate for a fixed teacher.

The finite-scale theorem proves a statement about clipped draws and
means of clipped outputs.  It makes no claim about the raw posterior
mean at finite scale.  For two or more hidden layers, this proof does not cover the entire
range $\gamma\geq1$: when $\gamma<2$, the quadratic parameter-count
term makes $P_N/N^\gamma$ grow with the learner width.  Removing that
term requires additional control in function space.  With one hidden
layer, that quadratic term is absent, and the theorem already covers
every $\gamma\geq1$.  The raw-output results
of Section~\ref{sec:raw-gaussian-oracle} instead concern the
strong-prior limit at each fixed width and therefore do not need
$\gamma\geq2$.

% Standalone sections for inclusion in the combined research paper.
% Requires amsmath, amssymb, amsthm and theorem/lemma/corollary environments.
% Bibliography key: hrw.

\section{A sharp oracle bound for the final Bayesian posterior mean}
\label{sec:alt-final-mean}

A different choice of prior and likelihood gives a substantially faster
approximation guarantee.  The predictor in this section is ordinary
Bayesian model averaging after all observations have been supplied.  It
uses neither a randomly selected training prefix nor an average of the
posteriors at successive times.  The guarantee is in expectation over the
training sample.

The underlying exponential-weights result is closely related to the
recent theorem of H\o gsgaard, Rebeschini and Wegel
\cite{hrw}, who prove optimal expected aggregation
rates for a finite dictionary.  We give a self-contained arbitrary-prior
version, with a conservative temperature constant, so that it applies
directly to continuous neural-network parameter distributions.  The
point of the application below is the explicit neural-network prior and
its dependence on teacher width, depth and weight magnitude.

Let $\pi$ be a probability measure on a measurable parameter space
$\Theta$, and let $(h,x)\mapsto f_h(x)$ be jointly measurable with values
in $[-1,1]$.  Let $D$ be any probability law of $(X,Y)$ with
$Y\in[-1,1]$.  For any bounded measurable predictor $f$, write
$R(f)=\mathbb E_D(f(X)-Y)^2$, and abbreviate $R(h)=R(f_h)$.
Define
\[
 R(h)=\mathbb E_D(f_h(X)-Y)^2,
 \qquad
 \widehat R_S(h)=\frac1n\sum_{i=1}^n(f_h(X_i)-Y_i)^2.
\]
For $n$ independent observations $S\sim D^n$, define
\begin{equation}
 p_S(dh)=
 \frac{\exp\{-\eta n\widehat R_S(h)\}\pi(dh)}
 {\int\exp\{-\eta n\widehat R_S(u)\}\pi(du)},
 \qquad
 \widehat f_S(x)=\int f_h(x)p_S(dh).
 \label{eq:alt-gibbs}
\end{equation}
For $n=0$, the posterior is understood to be $\pi$.

\begin{theorem}[Final posterior mean; arbitrary prior]
\label{thm:alt-final-mean}
Suppose $0<\eta\le \frac14\log(5/4)$.  For every integer $n\ge0$,
\begin{equation}
 \mathbb E_{S\sim D^n}R(\widehat f_S)
 \le
 \inf_Q\left\{
     \int R(h)Q(dh)
       +\frac{\mathrm{KL}(Q\Vert\pi)}{\eta(n+1)}
       \right\}.
 \label{eq:alt-final-oracle}
\end{equation}
The infimum is over probability measures $Q$ independent of $S$; they
may depend on $D$.  In particular, $\eta=1/20$ is admissible.
\end{theorem}

The leading coefficient of the comparison risk in
\eqref{eq:alt-final-oracle} is one.  The theorem allows arbitrary
misspecification and arbitrary dependence of the label on the input.
It concerns the posterior \emph{mean}; the same conclusion is not
asserted for a random posterior draw.  It is also an expectation bound,
not a high-probability bound.

\begin{lemma}[Deleting one observation]
\label{lem:alt-delete-one}
Let $q$ be any probability law on $[-1,1]$, let $y\in[-1,1]$, and set
\[
 p(da)=\frac{e^{-\eta(a-y)^2}q(da)}
                   {\int e^{-\eta(u-y)^2}q(du)}.
\]
If $4(e^{4\eta}-1)\le1$, then
\begin{equation}
 \left(\int a\,q(da)-y\right)^2
       \le\int(a-y)^2p(da).
 \label{eq:alt-delete-one}
\end{equation}
\end{lemma}

\begin{proof}
Write $\ell(a)=(a-y)^2\in[0,4]$, and interpolate from $p$ back to $q$ by
\[
 q_t(da)=\frac{e^{t\eta\ell(a)}p(da)}
                    {\int e^{t\eta\ell(u)}p(du)},
 \qquad 0\le t\le1.
\]
Thus $q_0=p$, $q_1=q$, and boundedness permits differentiation under the
integral:
\[
 \int\ell\,dq-\int\ell\,dp
       =\eta\int_0^1\operatorname{Var}_{q_t}(\ell)\,dt.
\]
The function $a\mapsto(a-y)^2$ is $4$-Lipschitz on $[-1,1]$.
The independent-pair formula for variance gives
$\operatorname{Var}_{q_t}(\ell)\le16\operatorname{Var}_{q_t}(a)$.
Furthermore,
\[
 \frac{dq_t}{dq}\le e^{4\eta(1-t)}.
\]
If probability measures $\nu,\mu$ satisfy $\nu\le K\mu$, then
$\operatorname{Var}_\nu(a)\le K\operatorname{Var}_\mu(a)$: use
$\operatorname{Var}_\nu(a)=\inf_b\int(a-b)^2d\nu$ and choose
$b=\int a\,d\mu$.  Consequently,
\[
 \int\ell\,dq-\int\ell\,dp
 \le16\eta\int_0^1e^{4\eta(1-t)}dt\,
                   \operatorname{Var}_{q}(a)
 =4(e^{4\eta}-1)\operatorname{Var}_{q}(a).
\]
Finally,
$(\int a\,dq-y)^2=\int\ell\,dq-\operatorname{Var}_q(a)$.
The assumed bound on $\eta$ proves the claim.
\end{proof}

\begin{proof}[Proof of Theorem~\ref{thm:alt-final-mean}]
Use $n+1$ independent observations, with full posterior $p$ and posterior
$q_i$ obtained by deleting observation $i$.  Conditional on all these
observations, the image of $p$ under $h\mapsto f_h(X_i)$ is the
exponential tilt in Lemma~\ref{lem:alt-delete-one} of the corresponding
image of $q_i$.  Hence, pointwise in the sample,
\[
 \frac1{n+1}\sum_{i=1}^{n+1}
     \biggl(\int f_h(X_i)q_i(dh)-Y_i\biggr)^2
 \le \int\widehat R_{S_{n+1}}(h)p(dh).
\]
The Gibbs variational principle gives, for every fixed comparison
measure $Q$ with finite relative entropy,
\[
 \int\widehat R_{S_{n+1}}\,dp
 \le \int\widehat R_{S_{n+1}}\,dQ
       +\frac{\mathrm{KL}(Q\Vert\pi)}{\eta(n+1)}.
\]
On taking expectations, each summand on the first display becomes
$\mathbb E_{S\sim D^n}R(\widehat f_S)$: $q_i$ uses $n$ independent
observations and is independent of observation $i$.  The fixed
comparison measure satisfies
$\mathbb E\int\widehat R_{S_{n+1}}dQ=\int R\,dQ$.
This proves the bound for every $Q$, and therefore for its infimum.
For the convenient constant, $e^{1/5}\le(1-1/5)^{-1}=5/4$, so
$\eta=1/20$ satisfies the hypothesis.
\end{proof}

\section{A continuous architecture prior with teacher-dependent rates}
\label{sec:alt-architecture}

We now specify a Bayesian neural-network learner for which the oracle
penalty depends on the teacher's parameters, not the largest architecture
available to the learner.  This is a prior on network architecture and
weights.  It differs from a fixed-width prior that independently draws
all weights from one Gaussian law.

\subsection{Network class and prior}

Inputs lie in any Borel subset $\mathcal X\subseteq[-1,1]^d$, where
$d\ge1$.  Fix a globally Lipschitz activation $\varphi$, and choose
$a\ge\max\{1,\operatorname{Lip}(\varphi)\}$ and
$b=|\varphi(0)|$.  No condition $\varphi(0)=0$ is needed.
A depth-$L$, width-$m$ network, with $L\ge2$ and $m\ge1$, has
$L-1$ hidden layers of width $m$ and a linear readout:
\[
 h^0(x)=(1,x),\qquad h^\ell(x)=\varphi(W^\ell h^{\ell-1}(x)),
 \quad 1\le\ell<L,\qquad z_w(x)=v^{\mathsf T}h^{L-1}(x).
\]
The activation is applied coordinatewise.  The input augmentation
provides first-layer biases; there are no other biases.  The parameter
count is
\begin{equation}
 P_{L,m}=m(d+1)+(L-2)m^2+m.
 \label{eq:alt-parameter-count}
\end{equation}
Its bounded prediction is $f_w(x)=\operatorname{clip}(z_w(x))$, where
$\operatorname{clip}(u)=\max\{-1,\min\{1,u\}\}$.

Let $N\in\mathbb N\cup\{\infty\}$ be the maximum allowed width, and
$L_{\max}\in\{2,3,\ldots\}\cup\{\infty\}$ the maximum allowed depth.  On
allowed pairs $(L,m)$ put prior mass proportional to
\begin{equation}
 \omega_{L,m}=\frac1{L(L-1)m(m+1)},
 \qquad 2\le L\le L_{\max},\quad 1\le m\le N.
 \label{eq:alt-architecture-weight}
\end{equation}
These weights sum to one when $N=L_{\max}=\infty$.  Conditioning on finite
architecture limits only increases the probability of each allowed
pair.  Given $(L,m)$, draw its $P_{L,m}$ weights independently from the
standard Cauchy law, with density $u\mapsto1/[\pi(1+u^2)]$.
Call the resulting proper prior $\pi_{N,L_{\max}}$.
The parameter space is the disjoint union of the allowed architecture
spaces.  In particular, varying $L$ selects networks of genuinely
different depths; it does not require embedding a shallower tanh
network into one fixed depth-$L_{\max}$ architecture.

The learner uses the likelihood
\begin{equation}
 \prod_{i=1}^n\exp\left\{-\frac1{20}
                  (f_w(X_i)-Y_i)^2\right\}
 \label{eq:alt-clipped-likelihood}
\end{equation}
and predicts with the final posterior mean of $f_w$.
This is the usual Gaussian likelihood for a clipped network mean and
observation variance $10$, up to factors independent of $w$.
The theorem below imposes no Gaussian assumption on the actual labels.
In particular it permits arbitrary label noise.  The learner does not
require exact interpolation of the observed labels.

\begin{theorem}[Agnostic neural-network oracle bound]
\label{thm:alt-nn-oracle}
Let $D$ be any probability law on $\mathcal X\times[-1,1]$, and let
$\widehat f_{n,N,L_{\max}}$ be the preceding posterior-mean predictor.  Let
$g=z_{w_*}$ be any network of depth $2\le L\le L_{\max}$ and width
$1\le m\le N$, all of whose weights have absolute value at most
$B\ge1$.  Put
\[
 D_{d,m}=\max\{d+1,m\},\qquad
 A=1+b+aD_{d,m}(B+1),\qquad K=L A^L,
\]
and let $P=P_{L,m}$ as in \eqref{eq:alt-parameter-count}.  For every
$0<\rho\le1$ and $n\ge0$,
\begin{align}
 \mathbb E R(\widehat f_{n,N,L_{\max}})
 &\le R(\operatorname{clip}g)+4\rho \\
 &\quad+\frac{20}{n+1}\left[
  \log\bigl(L(L-1)m(m+1)\bigr)
  +P\log\left(\frac{\pi[1+(B+1)^2]K}{2\rho}\right)
                              \right].
 \label{eq:alt-nn-bound}
\end{align}
The same inequality holds with $R(g)$ in place of
$R(\operatorname{clip}g)$.  Its right side is independent of $N$ and
$L_{\max}$ apart from the requirement that they admit the competitor.
\end{theorem}

The theorem holds for every choice of architecture limits, including
no finite limits at all.  Because it is an expectation bound, it does
not assert a common high-probability sample event for all choices of
these limits.  Taking the infimum in \eqref{eq:alt-nn-bound} permits a
different competitor at every sample size.

\begin{lemma}[Parameter perturbations]
\label{lem:alt-perturb}
In the setting of Theorem~\ref{thm:alt-nn-oracle}, if
$\|w-w_*\|_\infty\le r\le1$, then
\[
 \sup_{x\in\mathcal X}|z_w(x)-z_{w_*}(x)|\le K r.
\]
\end{lemma}

\begin{proof}
Both parameter vectors have coordinates bounded by $B+1$.
Let $M_\ell$ bound the maximum coordinate magnitude in hidden layer
$\ell$ of either network, and let $\delta_\ell$ bound their coordinatewise
difference, uniformly in $x$.  Since $M_0\le1$, the activation bound
$|\varphi(u)|\le b+a|u|$ gives $M_\ell\le A^\ell$ by induction.
The perturbation recurrence is
\[
 \delta_\ell\le aD_{d,m}r A^{\ell-1}
                    +aD_{d,m}(B+1)\delta_{\ell-1},
 \qquad \delta_0=0.
\]
Both $aD_{d,m}$ and $aD_{d,m}(B+1)$ are at most $A$, so induction gives
$\delta_\ell\le \ell r A^\ell$.
For the readout, use $D_{d,m}\le A$ and $D_{d,m}(B+1)\le A$ to obtain
\[
 |z_w-z_{w_*}|
 \le D_{d,m}r A^{L-1}+D_{d,m}(B+1)(L-1)r A^{L-1}
 \le LrA^L.
\]
\end{proof}

\begin{proof}[Proof of Theorem~\ref{thm:alt-nn-oracle}]
Choose a comparison measure supported on architecture $(L,m)$, uniform
on the parameter cube
$\{w:\|w-w_*\|_\infty\le r\}$ with $r=\rho/K\le1$.
Lemma~\ref{lem:alt-perturb} and the fact that clipping is nonexpansive
show $\|f_w-\operatorname{clip}g\|_\infty\le\rho$ throughout this cube.
Squared loss is $4$-Lipschitz in its prediction on $[-1,1]$, hence
\[
 \int R(f_w)Q(dw)\le R(\operatorname{clip}g)+4\rho.
\]
On the cube, each Cauchy coordinate density is at least
$1/\{\pi[1+(B+1)^2]\}$.  The uniform density there is $(2r)^{-P}$.
Moreover, the architecture has prior mass at least $\omega_{L,m}$.
Therefore
\[
 \mathrm{KL}(Q\Vert\pi_{N,L_{\max}})
 \le \log\frac1{\omega_{L,m}}
       +P\log\left(\frac{\pi[1+(B+1)^2]}{2r}\right).
\]
Apply Theorem~\ref{thm:alt-final-mean} with $\eta=1/20$ and substitute
$r=\rho/K$.  Finally, clipping cannot increase squared error against a
label in $[-1,1]$, so $R(\operatorname{clip}g)\le R(g)$.
\end{proof}

\subsection{The same rate with conditionally Gaussian weights}

The logarithmic cost of the teacher's weight magnitude does not require
Cauchy coordinates.  There is also a Gaussian scale-mixture version of
the prior.  After selecting $(L,m)$ with the same architecture law,
choose an integer $j\ge0$ with probability $1/[(j+1)(j+2)]$, and then
draw the weights independently from $N(0,4^j)$.  The prior is proper,
and its coordinates are independent conditional on the architecture
and scale.

\begin{corollary}[Gaussian scale mixture]
\label{cor:alt-gaussian-mixture}
Use this Gaussian scale-mixture prior with the likelihood and predictor
of \eqref{eq:alt-clipped-likelihood}.  For the competitor in
Theorem~\ref{thm:alt-nn-oracle}, let
$j_B=\lceil\log_2 B\rceil$.  Then
\begin{align}
 \mathbb E R(\widehat f_{n,N,L_{\max}})
 &\le R(\operatorname{clip}g)+4\rho\\
 &\quad+\frac{20}{n+1}\left[
   \log(1/\omega_{L,m})
   +\log\bigl((j_B+1)(j_B+2)\bigr)
   +P\log\left(\frac{B\sqrt{2\pi}e^2K}{\rho}\right)
                           \right].
 \label{eq:alt-gaussian-mixture}
\end{align}
\end{corollary}

\begin{proof}
The comparison measure selects architecture $(L,m)$, selects scale
$j_B$, and uses the same uniform parameter cube of radius $r=\rho/K$.
Put $s=2^{j_B}$, so $B\le s\le2B$.
Every coordinate in the cube satisfies $|w_i|\le B+1\le2s$; its
$N(0,s^2)$ density is therefore at least
$e^{-2}/(s\sqrt{2\pi})$.  The coordinate contribution to relative
entropy is at most
\[
 P\log\left(\frac{s\sqrt{2\pi}e^2}{2r}\right)
 \le P\log\left(\frac{B\sqrt{2\pi}e^2}{r}\right).
\]
Adding the architecture and scale selection costs gives the bracket
in \eqref{eq:alt-gaussian-mixture}.  The risk comparison and final
posterior-mean theorem apply as before.
\end{proof}

The additional scale cost is only
$O(1+\log(1+\log B))$.  Both prior choices consequently have the
same polynomial width--depth rates described next.  A single fixed
Gaussian scale would instead charge order $PB^2$ for translating the
comparison cube; the scale mixture makes this cost logarithmic in $B$.

\subsection{Rates, activations, and what changes in the learner}

Setting $\rho=1/(n+1)$ in \eqref{eq:alt-nn-bound} yields a comparison
penalty of order
\begin{equation}
 \frac{1}{n+1}\left\{
 1+\log(Lm)
 +P_{L,m}\left[
   \log(n+1)+\log(B+1)+\log L
      +L\log A
                  \right]\right\}.
 \label{eq:alt-rate-summary}
\end{equation}
Here $A=1+b+a\max\{d+1,m\}(B+1)$, as above, and the hidden
constant is numerical.  Thus the cost is approximately one
bit per parameter per required precision bit, together with the
architecture description.  The dependence on the weight bound is
logarithmic.  No bound on $B$ must be selected in advance by the learner.

For fixed activation, input dimension and depth, competitors of width
$m\le n^\alpha$ and weights bounded by $n^\beta$, where $\beta$ is any
fixed nonnegative constant, have excess expected risk
$O(n^{2\alpha-1}\log n)$.  Every $\alpha<1/2$ is therefore admissible.
For a single hidden layer ($L=2$),
$P_{2,m}=m(d+2)$, and this improves to
$O(n^{\alpha-1}\log n)$: every $\alpha<1$ is admissible.
If the depth also varies, the displayed bound has polynomial, rather
than exponential, depth dependence.  Suppressing logarithms and holding
$d$, $B$ and the activation fixed, it is
$\widetilde O((mdL+m^2L^2)/(n+1))$.
These statements are approximation guarantees: the target distribution
need not be representable by any network.

The activation assumption includes tanh and ReLU with $a=1,b=0$;
GELU, $\varphi(u)=u\Phi(u)$, with $a=2,b=0$; and softplus with
$a=1,b=\log2$.  In particular, the proof neither relies on bounded
hidden activations nor requires $\varphi(0)=0$.  The Gaussian derivative
bound $|\Phi(u)+u\phi(u)|\le1+(2\pi e)^{-1/2}<2$ verifies the stated
GELU constant.  Internal biases can be added by augmenting each layer
with a constant coordinate and changing the parameter count and
$D_{d,m}$ accordingly.

The price of this improved bound is a change of prior and prediction
convention, not an unproved property of the original dense Gaussian
prior.  A draw from $\pi_{N,L_{\max}}$ first selects the \emph{active}
architecture and then draws its weights.  At a fixed depth and maximum width it
can equivalently be regarded as a structural sparsity prior, with exact
zero connections outside the selected subnetwork.  It is not a product
law on all the weights of a dense width-$N$ network.  Its preference for
small architectures is explicit.  The guarantee therefore does not
resolve whether an unrestricted dense Gaussian prior has a comparable
polynomial-depth approximation theorem.

Finally, the loss in the likelihood is computed from the clipped
network output, and the predictor averages those clipped outputs.
Replacing either with an unrestricted real output is a different
learner.  Boundedness is used in Lemma~\ref{lem:alt-delete-one}; no such
replacement follows from that lemma.

\section{What these guarantees establish}

\subsection{Approximation rather than exact realization}
The comparison network need not describe the entire data-generating process. It may simply be a useful approximation. For example, at fixed depth and dimension, the architectural-prior theorem gives
\[
 \mathbb E R(\widehat f_n)
 \le \inf_{g}\left\{R(\operatorname{clip}g)
       +\widetilde O\!\left(\frac{m(g)^2}{n}\right)\right\},
\]
when the comparison weights grow at most polynomially with $n$; the logarithms and depth factors are specified in Theorem~\ref{thm:alt-nn-oracle}. Increasing the teacher width can reduce approximation error while increasing the estimation penalty. This is the usual approximation--estimation tradeoff, now stated for a particular Bayesian learner. With a single hidden layer the number of parameters is linear in width, so the corresponding penalty is linear in $m$ at fixed dimension.

The Gaussian theorem obtains the same kind of comparison without introducing a prior over architectures, but with a slower depth-dependent rate. The two guarantees therefore answer different versions of the learning question. Neither supplies an efficient posterior-sampling algorithm or an optimization guarantee for gradient descent.

\subsection{Why the likelihood is not an interpolation requirement}
With arbitrary label noise, fitting every observation closely is a different objective from predicting its conditional mean. For the Gaussian learner, the likelihood variance is $n\lambda_n\sigma^2$. Although this tends to zero as $\sigma\downarrow0$, the prior simultaneously strengthens; the posterior selects minimizers of empirical squared loss plus $\lambda_n Q_N$, rather than exact interpolants. The regularization term is essential to our proof.

For the architectural prior we instead use a fixed observation variance of $10$, after normalizing labels and predictions to $[-1,1]$. This is an ordinary Gaussian observation likelihood used under possible misspecification. Its relatively weak per-observation weighting permits the posterior-mean averaging argument. Neither construction assumes that the actual data noise is Gaussian or has the likelihood's variance.

Earlier counterexamples~\cite{opus} show that, at depth at least three and sufficiently large learner width, the untempered strong-prior objective can memorize independent label noise and have population error bounded away from that of the zero predictor. The positive results here do not turn exact noisy interpolation into a generally valid agnostic learning rule.

\subsection{A width-only comparison needs a qualification on weights}
Our Gaussian bound charges teacher weight magnitude polynomially. The architectural-prior bound charges it only logarithmically. Some dependence on precision or weight magnitude is necessary for a distribution-uniform guarantee for a fixed proper Bayesian prior on continuous functions.

\begin{proposition}[An obstruction from nearby inputs]
\label{prop:nearby-inputs}
Fix $n<\infty$, $\eta>0$, and any proper prior $\pi$ on continuous predictors $f:\mathbb R\to[-1,1]$, with measurable evaluation. Use the posterior proportional to
$\exp[-\eta\sum_{i=1}^n(f(X_i)-Y_i)^2]\pi(df)$.
Let $D_r$ assign probability $1/2$ to each of $(-r,-1)$ and $(r,1)$.
As $r\downarrow0$, uniformly over every sample of size $n$ from $D_r$, the posterior-mean risk and the posterior-averaged risk satisfy
\[
 \liminf R_{D_r}(\widehat f_S)\ge1,
 \qquad
 \liminf \int R_{D_r}(f)\,p_S(df)\ge1.
\]
Nevertheless a width-one tanh network at any prescribed fixed depth has zero risk under $D_r$, if its weights may grow as $r\downarrow0$.
\end{proposition}
\begin{proof}
The likelihood lies in $[e^{-4\eta n},1]$, so $dp_S/d\pi\le e^{4\eta n}$, uniformly over the sample. Continuity and bounded convergence give
\[
 a_r:=\int|f(r)-f(-r)|\,\pi(df)\longrightarrow0.
\]
For every bounded predictor,
\[
 R_{D_r}(f)=1+\tfrac12\bigl(f(-r)^2+f(r)^2\bigr)+f(-r)-f(r)
 \ge1-|f(r)-f(-r)|.
\]
Integrating proves the posterior-draw bound $1-e^{4\eta n}a_r$. The posterior mean satisfies the same bound by the triangle inequality. For the teacher, set the first scalar weight to $1/r$, the constant-input weight to zero, and all remaining hidden weights to one. If $T$ denotes tanh composed $L-1$ times, choose readout weight $1/T(1)$. Its two outputs are exactly $-1$ and $1$. Its first weight grows like $1/r$.
\end{proof}

The obstruction applies to the architectural priors above as well: every sampled finite network is continuous, even when the prior has no maximum width or depth. It does not contradict their oracle bounds, whose penalty depends on the large teacher weight. It is not an impossibility result for every learning algorithm; it concerns the stated fixed-prior Bayesian rule. The same obstruction holds on the unit circle: use the inputs $(\pm r,\sqrt{1-r^2})$, which coalesce, and the same teacher applied to the first coordinate.

\subsection{Expectation, high probability, and the remaining Gaussian question}
The Gaussian result provides a high-probability bound on one sample event for all eligible learner widths. The architectural-prior result instead bounds the expected risk of each learner. These quantifiers cannot be interchanged. In particular, the latter proof does not give the same rate with only a $\log(1/\delta)$ confidence penalty, and its posterior-draw risk need not obey the same bound. The posterior mean can have lower risk than every individual member of its support; its improvement is used explicitly in the leave-one-out proof.

For dense independent Gaussian priors, the principal remaining question is whether one can obtain a comparable teacher-parameter-count oracle bound without explicitly selecting a small architecture. The present paper improves the strong-prior rate, but does not prove such a fast rate or a joint polynomial-depth guarantee for that learner. The architectural result is not a resolution of this Gaussian-prior question.

\paragraph{Verification and provenance.}
The positive theorem in the supplied tempered-learner note was checked against its complete supporting proof~\cite{opus}; no error was found in that theorem under its stated conventions. The stronger results here have independent AI-agent mathematical reviews and reproducible source builds. They have not been formalized in Lean and have not undergone external peer review. The supporting audit records distinguish this scope from the collection's formally verified results.


\appendix
\section{The same sample exponent without clipping}
\label{sec:raw-gaussian-oracle}

The Gaussian-prior model and likelihood family in Section~\ref{sec:gaussian-oracle}, with the calibration below, also admit
an oracle inequality for its \emph{raw} real-valued predictions.  This
covers both posterior-averaged risk and the ordinary posterior mean
$\int z_w\,d\Pi_{\sigma,N,S}$, with no clipping in either the likelihood
or the predictor.  The improvement uses a standard uniform bound for
smooth losses.  It has the same sample exponent as the elementary
clipped argument, at the cost of additional logarithms.

We retain the architecture, energy $Q_N$, prior, likelihood, and risks
from Section~\ref{sec:gaussian-oracle}.  There are $H$ hidden layers and
one linear readout.  All labels lie in $[-1,1]$, and the joint law of
inputs and labels is otherwise arbitrary.  A teacher $g$ of width $m$
and weights bounded by $B$ has clones of energy at most
$E=d(d+1)B^2+2Hm^2B^2$ at every learner width $N\geq m$, and satisfies
$|g|\leq U=mB$.

\subsection{The smooth-loss inequality}

We use the following form of Theorem~1 of
Srebro, Sridharan, and Tewari~\cite[Theorem~1, p.~4 of the full-length version]{sst}.  If a nonnegative loss is
$s$-smooth in its real-valued prediction and is bounded by $b$ on a
function class $\mathcal F$, then, with probability at least $1-\eta$,
simultaneously for $f\in\mathcal F$,
\begin{align}
 R(f)-\widehat R_S(f)
 &\leq C_0\biggl[
  \sqrt{\widehat R_S(f)}
    \left(\sqrt{s}\log^{3/2}(n)\mathfrak R_n(\mathcal F)
           +\sqrt{\frac{b\log(1/\eta)}n}\right)
 \notag\\[-2pt]
 &\hspace{37mm}
  +s\log^3(n)\mathfrak R_n(\mathcal F)^2
  +\frac{b\log(1/\eta)}n\biggr].
 \label{eq:raw-sst}
\end{align}
Here $\mathfrak R_n$ is the worst-case Rademacher complexity and $C_0$ is
universal; $C_0=10^5$ is sufficient.  This result assumes neither a convex
hypothesis class nor a realizable data law.  In particular, it does not
assume a Bernstein condition on excess loss relative to the teacher.

The raw square loss is nonnegative and $2$-smooth.  On $\mathcal Z_K$
from~\eqref{eq:gaussian-Dn}, its range is bounded by
$(\sqrt K+1)^2$, since $|z_w|\leq\sqrt{Q_N(w)}$.
The class $\mathcal Z_K$ is symmetric under readout sign reversal, so
its Rademacher complexity with an absolute value inside the supremum,
as used in~\cite{sst}, equals the convention used in
Lemma~\ref{lem:gaussian-capacity}.

\subsection{The theorem for two or more hidden layers}

For $n\geq2$, $0<\delta<1$, and $H\geq2$, define
\begin{gather}
 p=H/2,\qquad J_n=\lceil\log_2 n\rceil,\qquad
 t_n=\log\frac{2(J_n+1)}\delta,\qquad q_n=t_n/n,\notag\\
 A_n=\frac{\sqrt2\,D_n\log^{3/2}(2n)}{H^p\sqrt n},\qquad
 a_n=2^p A_n,\notag\\
 r_n=C_0\bigl(A_n+2\sqrt{q_n}+A_n^2+4q_n\bigr).
 \label{eq:raw-constants}
\end{gather}
For a fixed teacher $g$, put
\begin{equation}
 b_n(g,\delta)
  =(U+1)\sqrt{\frac{2R(g)\log(2/\delta)}n}
     +\frac{2(U+1)^2\log(2/\delta)}{3n}.
 \label{eq:raw-teacher-concentration}
\end{equation}
The slightly conservative coefficient $2/3$ makes this a direct
consequence of the usual Bernstein tail inequality.

\begin{theorem}[A raw-output oracle inequality at every learner width]
\label{thm:raw-gaussian-oracle}
Let $H\geq2$, and suppose the likelihood parameter $\lambda$ satisfies
\begin{equation}
 \frac1n\leq\lambda\leq1,\qquad
 \lambda^p\geq8C_0a_n,\qquad \lambda\geq32C_0q_n.
 \label{eq:raw-calibration}
\end{equation}
For every fixed teacher $g$ as above, with probability at least
$1-\delta$ over the sample, simultaneously for every $N\geq m$,
\begin{equation}
 \limsup_{\sigma\downarrow0}
   \int R(z_w)\,\Pi_{\sigma,N,S}(dw)
 \leq R(g)+\lambda E+r_n+\frac{8C_0^2q_n}{\lambda}
                 +b_n(g,\delta).
 \label{eq:raw-oracle}
\end{equation}
The same right-hand side bounds the limiting risk of the ordinary
posterior mean $\int z_w\,d\Pi_{\sigma,N,S}$.  It also bounds the limiting
risk after clipping that mean, and the limiting risk of the posterior
mean of clipped outputs.

An explicit choice is
\begin{equation}
 \lambda_n=\max\bigl\{(8C_0a_n)^{1/p},\ 32C_0q_n\bigr\},
 \label{eq:raw-explicit-schedule}
\end{equation}
for all sufficiently large $n$ that $\lambda_n\leq1$.  At fixed
$d,H,\delta$, this gives an excess-risk bound of order
$\widetilde O(n^{-1/H})$ for every fixed teacher.  More explicitly,
\[
 \lambda_n
  =O_H\!\left[
      \left(\frac{D_n^2\log^3(2n)}n\right)^{1/H}
      +\frac{\log(\log(2n)/\delta)}n\right].
\]
The tilde suppresses logarithms; it does not suppress a learner-width
factor.  There is no dependence on $N$ in the displayed bound.
\end{theorem}

\begin{proof}
Apply~\eqref{eq:raw-sst} to each deterministic energy class
$\mathcal Z_{2^j}$, $0\leq j\leq J_n$, using failure probability
$\delta/[2(J_n+1)]$.  A union bound gives one event of probability at
least $1-\delta/2$ for all these classes and all learner widths.

Let $w$ minimize $\widehat R_S(z_w)+\lambda Q_N(w)$.  Comparing with
the zero network gives
\begin{equation}
 \widehat R_S(z_w)+\lambda Q_N(w)
   \leq\widehat R_S(0)\leq1.
 \label{eq:raw-training-budget}
\end{equation}
Write $k=Q_N(w)$.  If $k\leq1$, the event for $\mathcal Z_1$ implies
$R(z_w)-\widehat R_S(z_w)\leq r_n$.
If $k>1$, choose $j$ with $2^{j-1}<k\leq2^j$.  This is possible because
$k\leq1/\lambda\leq n$.  The shell radius $K=2^j$ satisfies $K\leq2k$.
The loss bound obeys $(\sqrt K+1)^2\leq4K\leq8k$.
Substituting Lemma~\ref{lem:gaussian-capacity} and
$\widehat R_S(z_w)\leq1$ into~\eqref{eq:raw-sst} yields
\begin{equation}
 R(z_w)-\widehat R_S(z_w)
 \leq C_0\left[a_n k^p+a_n^2 k^{2p}
                    +2\sqrt2\sqrt{kq_n}+8kq_n\right].
 \label{eq:raw-shell-bound}
\end{equation}
Since $p\geq1$, $k\leq1/\lambda$, and
$\lambda^p\geq8C_0a_n$, the first two terms satisfy
\[
 C_0a_n k^p\leq\frac{\lambda k}8,
 \qquad
 C_0a_n^2 k^{2p}\leq\frac{\lambda k}{64C_0}
                     \leq\frac{\lambda k}{64}.
\]
The remaining terms obey
\[
 8C_0kq_n\leq\frac{\lambda k}4,
 \qquad
 2\sqrt2 C_0\sqrt{kq_n}
  \leq\frac{\lambda k}4+\frac{8C_0^2q_n}{\lambda}.
\]
Consequently both the $k\leq1$ and $k>1$ cases imply
\begin{equation}
 R(z_w)\leq\widehat R_S(z_w)+\lambda Q_N(w)
                      +r_n+\frac{8C_0^2q_n}{\lambda}.
 \label{eq:raw-absorbed-deviation}
\end{equation}

Compare the penalized empirical objective with the exact teacher clone:
\[
 \widehat R_S(z_w)+\lambda Q_N(w)
       \leq\widehat R_S(g)+\lambda E.
\]
The teacher's loss lies in $[0,(U+1)^2]$ and has variance at most
$(U+1)^2R(g)$.  Bernstein's inequality therefore gives
$\widehat R_S(g)\leq R(g)+b_n(g,\delta)$ with probability at least
$1-\delta/2$.  Intersecting the two events proves the right-hand side
of~\eqref{eq:raw-oracle} for every global minimizer, at every width
$N\geq m$.

For completeness, the passage to the posterior also holds for raw
risk.  At fixed $n,N,S$, the objective is continuous and coercive.
Its minimizer set is compact, and its near-minimum sublevel sets shrink
to that set.  The risk $w\mapsto R(z_w)$ is continuous on every compact
parameter set by dominated convergence, using
$|z_w|\leq\sqrt{Q_N(w)}$.  Moreover,
\[
 R(z_w)\leq2Q_N(w)+2.
\]
To make the required tail control explicit, write
$F(w)=Q_N(w)+\lambda^{-1}\widehat R_S(z_w)$ and
$F_* = \min_w F(w)$.  For every $\eta>0$,
\begin{equation}
 \lim_{\sigma\downarrow0}
 \int (1+Q_N(w))\,
       \mathbf 1_{\{F(w)>F_*+\eta\}}\,
       \Pi_{\sigma,N,S}(dw)=0.
 \label{eq:raw-weighted-laplace-tail}
\end{equation}
Indeed, a fixed ball around a minimizer lower-bounds the normalizer by
its volume times $\exp[-(F_*+\eta/4)/(2\sigma^2)]$.
Put $\rho=\eta/[4(F_*+\eta)]$.  On the indicated tail,
$F\geq(1-\rho)(F_*+\eta)+\rho Q_N
      =F_*+3\eta/4+\rho Q_N$.
The remaining integral is a quadratic Gaussian integral of $1+Q_N$,
which grows or decays at most polynomially in $\sigma$ and is dominated
by the factor $\exp[-\eta/(4\sigma^2)]$ after division by the
normalizer.  This proves~\eqref{eq:raw-weighted-laplace-tail}.
Consequently the limiting posterior average of the risk is at most its
maximum over the minimizer set, which proves~\eqref{eq:raw-oracle}.  Jensen's inequality and Tonelli's
theorem give the conclusion for the ordinary posterior mean.  Clipping
can only decrease squared loss for labels in $[-1,1]$; applying the same
facts gives the other prediction conventions.

Finally,~\eqref{eq:raw-explicit-schedule} satisfies
$\lambda_n\geq1/n$ and the other lower bounds in
\eqref{eq:raw-calibration}.  Its leading power is $n^{-1/H}$.
The terms $r_n$ have powers at least $n^{-1/2}$, up to logarithms, and
$q_n/\lambda_n$ has power at least $n^{-1+1/H}$.  Both are no slower than
$n^{-1/H}$ when $H\geq2$.
\end{proof}

The teacher in this theorem may be selected using the data law and
$n$, but must be fixed before the sample is drawn.  In particular, it
may minimize the population risk over a specified compact class of
bounded-weight teachers.  Such a class contains the zero network, so
its minimizer has risk at most one.  This interpretation gives a sharp
oracle inequality with leading coefficient one, rather than a
realizability guarantee.  The event is simultaneous over learner widths;
no simultaneity over all teachers is needed for this interpretation.

\subsection{One hidden layer}

The same smooth-loss argument gives the rate $\widetilde O(n^{-1/2})$
for raw predictions when $H=1$.  Define $p=1/2$ and the quantities
$A_n,a_n,q_n,r_n$ by~\eqref{eq:raw-constants}.  The following explicit
version records the different absorption step.

\begin{theorem}[Raw predictions with one hidden layer]
\label{thm:raw-gaussian-one-hidden}
Let $H=1$ and suppose
\[
 1/n\leq\lambda\leq1,
 \qquad \lambda\geq2C_0(a_n^2+8q_n).
\]
For every fixed teacher $g$, with probability at least $1-\delta$,
simultaneously for every $N\geq m$,
\begin{align}
 \limsup_{\sigma\downarrow0}\int R(z_w)\,\Pi_{\sigma,N,S}(dw)
 \leq R(g)+\lambda E+r_n
       +\frac{C_0^2(a_n+2\sqrt{2q_n})^2}{2\lambda}
       +b_n(g,\delta).
 \label{eq:raw-one-hidden-oracle}
\end{align}
The ordinary posterior mean and the clipped prediction conventions
satisfy the same bound.  For example,
\[
 \lambda_n=32C_0(a_n+\sqrt{q_n}+a_n^2+q_n)
\]
satisfies the lower bounds and gives the rate
$\widetilde O(n^{-1/2})$ for a fixed teacher, whenever $n$ is sufficiently
large that $\lambda_n\leq1$.
\end{theorem}
\begin{proof}
The same simultaneous shell event yields~\eqref{eq:raw-shell-bound}.
For $p=1/2$, its right-hand side is
\[
 C_0(a_n+2\sqrt{2q_n})\sqrt k+C_0(a_n^2+8q_n)k.
\]
The second term is at most $\lambda k/2$.  Young's inequality bounds
the first by
\[
 \frac{\lambda k}2+
       \frac{C_0^2(a_n+2\sqrt{2q_n})^2}{2\lambda}.
\]
The rest is the teacher comparison, scalar concentration, and Laplace
argument from the preceding proof.  For the stated schedule,
$a_n=\widetilde O(n^{-1/2})$ and
$q_n=\widetilde O(n^{-1})$, giving the claimed rate.
\end{proof}

If a teacher-energy upper bound $E_0\geq E$ is specified in advance,
there is a useful tuning that improves the teacher dependence.  Put
\begin{gather*}
 V_n=\frac{C_0^2(a_n+2\sqrt{2q_n})^2}{2},\qquad
 v_n=\max\{2C_0(a_n^2+8q_n),\,1/n\},\\
 \lambda_n=\max\left\{\sqrt{\frac{V_n}{E_0+1}},\,v_n\right\}.
\end{gather*}
Whenever $\lambda_n\leq1$, Theorem~\ref{thm:raw-gaussian-one-hidden}
applies, and
\[
 \lambda_n E+\frac{V_n}{\lambda_n}
   \leq2\sqrt{V_n(E_0+1)}+v_nE.
\]
At fixed input dimension, this contribution is
$\widetilde O(\sqrt{(E_0+1)/n}+E_0/n)$.
For example, for population-optimal teachers whose width and weights
are both at most $n^\alpha$, one may take $E_0=O_d(n^{4\alpha})$.
The teacher risk is then at most one, and all terms in
\eqref{eq:raw-one-hidden-oracle} tend to zero for every $\alpha<1/4$,
with leading power $n^{2\alpha-1/2}$.


\begin{thebibliography}{9}
\bibitem{opus}
Claude Opus 5.5.
\emph{Tempered-likelihood agnostic learning for Bayesian tanh networks}.
Research note and short theorem statement, November 2026.
Complete supporting proof: \href{https://github.com/kaarelh/AI-works/blob/main/neural-network-generalization/works-in-progress/agnostic-learning/research-notes/B-tempered-likelihood.md}{\texttt{B-tempered-likelihood.md}}, Sections 2--6 and 8.

\bibitem{hrw}
Mikael M\o ller H\o gsgaard, Patrick Rebeschini and Tobias Wegel.
\emph{Aggregation with Exponential Weights is Optimal in Expectation}.
2026. \href{https://arxiv.org/abs/2607.02247}{arXiv:2607.02247}.

\bibitem{sst}
Nathan Srebro, Karthik Sridharan and Ambuj Tewari.
\emph{Smoothness, Low Noise and Fast Rates}.
Advances in Neural Information Processing Systems 23, 2010.
\url{https://www.cs.cornell.edu/~sridharan/smoothness.pdf}.
\end{thebibliography}
\end{document}
